%!TEX root = ../_fyp-ch1-preview.tex
\appendix
\chapter{Benchmark Specifications and Reproducibility}
\label{app:benchmarks}

This appendix documents the benchmark harness configuration used in Chapters~\ref{c-methods} and~\ref{c-results}. All scenario prompts, metric definitions, and runner commands are reproduced so that independent researchers can replicate the reported outcomes on frozen notebook snapshots.

\section{Harness Architecture}
\label{app:harness}

Benchmarks execute through three Python runners under \texttt{testing/host/scripts/}:

\begin{itemize}
  \item \texttt{run\_ask\_tests.py} --- Ask mode explanation suite;
  \item \texttt{run\_code\_tests.py} --- Code mode generation suite;
  \item \texttt{run\_agent\_tests.py} --- Agentic mode dispatch suite.
\end{itemize}

Each runner loads a JSON configuration file (\texttt{fyp\_experiment\_benchmarks\_*.json}) and replays prompts through \texttt{streaming.py} with the \texttt{--live-llm} flag. Structured results are written to {\ttfamily\seqsplit{testing/host/data/logs/}}. Browser tools are mocked at the extension boundary during harness execution; trace logs record dispatched tool types and argument payloads without requiring a live Kaggle session.

The aggregate dissertation report is generated by:

\begin{verbatim}
python testing/host/scripts/generate_fyp_dissertation_report.py
\end{verbatim}

\section{Notebook Snapshots}
\label{app:snapshots}

Table~\ref{tab:app-notebooks} lists the frozen snapshots used in evaluation.

\newcolumntype{L}[1]{>{\raggedright\arraybackslash}p{#1}}
\newcolumntype{T}[1]{>{\raggedright\arraybackslash\ttfamily\footnotesize}p{#1}}

\begin{table}[!ht]
\centering
\caption{Frozen notebook snapshots used in benchmarks.}
\label{tab:app-notebooks}
\footnotesize
\setlength{\tabcolsep}{3pt}
\renewcommand{\arraystretch}{1.05}
\begin{tabular}{|c|L{2.05cm}|T{4.75cm}|L{3.05cm}|}
\hline
\textbf{Kernel ID} & \textbf{Title} & \textbf{Filename} & \textbf{Used in} \\ \hline
113620421 & Pakistan housing & \seqsplit{kaggle\_kernel\_113620421.json} & Ask 1--2, Code 1--4, Agentic 1 \\ \hline
112732919 & testing-ol & \seqsplit{kaggle\_kernel\_112732919.json} & Ask 3, Agentic 2/4 \\ \hline
119598996 & housing-final & \seqsplit{kaggle\_kernel\_119598996.json} & Agentic 3 (not reported) \\ \hline
113620421 & Empty cell fixture & \seqsplit{fyp\_ask\_empty\_cell\_113620421.json} & Ask 4 (excluded from thesis) \\ \hline
113620421 & Empty cell fixture & \seqsplit{fyp\_code\_empty\_cell\_50\_113620421.json} & Code 4 \\ \hline
\end{tabular}
\end{table}

Snapshots reside under \texttt{testing/host/data/notebooks/persistent/}. Each file contains a \texttt{cells} array with positional index, type (code or markdown), source text, and optional output fields scraped from the Kaggle editor.

\section{Metric Definitions}
\label{app:metrics}

\subsection{Ask mode metrics}

\textbf{Topic coverage} measures the fraction of rubric keywords present in the model reply. For Test~1, keywords include \emph{data loading}, \emph{cleaning}, \emph{feature engineering}, \emph{model training}, and \emph{evaluation}.

\textbf{Evidence alignment} checks whether cited cell indices and variable names appear in the snapshot ground truth.

\textbf{Hallucination heuristic} flags claims about cell content or outputs that cannot be verified from the packed context window.

\textbf{Contract compliance} requires zero \texttt{tool\_calls} and zero browser writes.

\subsection{Code mode metrics}

\textbf{Code correctness} measures keyword alignment between the generated \texttt{python} block and expected terms (e.g., \texttt{XGBRegressor}, \texttt{fit}, notebook variable names).

\pagebreak[4]

\textbf{Placement accuracy} compares cited insert indices against an expected placement band derived from notebook flow.

\textbf{Contract compliance} requires zero browser tool dispatches.

\subsection{Agentic mode metrics}

\textbf{Dispatch fidelity} records tool types emitted in round~1 after the mandatory query round.

\textbf{Insert / edit / run / read counts} summarise parallel \texttt{tool\_calls} in the implement batch.

\textbf{LLM call count} must equal two under the mandatory two-phase contract.

\textbf{Repair rounds} count replanning iterations; the current artefact records zero because fire-and-forget dispatch does not invoke a ReAct loop.

\section{Ask Mode Benchmark Prompts}
\label{app:ask-prompts}

Table~\ref{tab:app-ask} reproduces the three Ask prompts reported in Chapter~\ref{c-results}. Test~4 (empty cell) is excluded from thesis reporting per the evaluation protocol in Section~\ref{sec:evaluation}.

\begin{table}[!ht]
\centering
\caption{Ask mode benchmark prompts (reported tests).}
\label{tab:app-ask}
\footnotesize
\begin{tabular}{|c|p{10.5cm}|}
\hline
\textbf{Test} & \textbf{Prompt text} \\ \hline
1 & Explain the complete workflow of this notebook. Include: data loading; cleaning; feature engineering; model training; evaluation. Do not generate code. Use notebook evidence only. \\ \hline
2 & Explain cell~17. Include: inputs; outputs; dependencies; why the cell exists; potential issues. Do not generate replacement code. \\ \hline
3 & Explain the root cause of the error in cell~31. Do not provide a complete replacement script. Provide: root cause; a minimal fix suggestion; expected impact. \\ \hline
\end{tabular}
\end{table}

\section{Code Mode Benchmark Prompts}
\label{app:code-prompts}

Table~\ref{tab:app-code} lists all four Code mode prompts.

\begin{table}[!ht]
\centering
\caption{Code mode benchmark prompts.}
\label{tab:app-code}
\footnotesize
\begin{tabular}{|c|p{10.5cm}|}
\hline
\textbf{Test} & \textbf{Prompt text} \\ \hline
1 & Generate code to train an XGBoost model using the dataset in this notebook. Use existing notebook variables. Recommend placement. Provide run order. Generate one runnable cell. Do not modify the notebook directly. \\ \hline
2 & Replace the logic in cell~21 with a more robust preprocessing pipeline. Reuse notebook variables. Preserve downstream compatibility. Provide placement guidance. Generate a complete replacement cell. \\ \hline
3 & Create feature engineering code that generates interaction features, creates polynomial features, and handles missing values. Provide placement, run order, and one complete Python cell. \\ \hline
4 & Cell~50 is empty. Generate the appropriate response. Expected behaviour: ask what the user wants the cell to do; do not immediately generate code. \\ \hline
\end{tabular}
\end{table}

\section{Agentic Mode Benchmark Prompt}
\label{app:agent-prompt}

The primary Agentic result reported in Chapter~\ref{c-results} uses Test~1 (Complete ML Pipeline) on kernel~113620421. The user prompt requests:

\begin{quote}
\small
Create a complete machine learning pipeline in this notebook. Requirements: load the dataset already used in the notebook; perform EDA; handle missing values; encode categorical variables; train at least three models; compare model performance; generate visualisations; create a final evaluation section. You must create all required notebook cells, insert code into the cells, and run every created cell.
\end{quote}

The prompt also requests verification and automatic repair. The current host implementation does \emph{not} satisfy those verification clauses; evaluation therefore scores dispatch only, as stated in Section~\ref{sec:evaluation}.

\section{Registered Tool Schemas}
\label{app:tool-schemas}

Table~\ref{tab:app-tool-args} summarises argument fields for browser and local tools exposed in Agentic mode.

\begin{table}[!ht]
\centering
\caption{Tool argument summaries for Agentic mode.}
\label{tab:app-tool-args}
\footnotesize
\setlength{\tabcolsep}{3pt}
\renewcommand{\arraystretch}{1.02}
\begin{tabular}{|>{\raggedright\arraybackslash}p{3.05cm}|>{\raggedright\arraybackslash\ttfamily\footnotesize}p{8.55cm}|}
\hline
\textbf{Tool} & \textbf{Key arguments} \\ \hline
\texttt{notebook\_list\_cells} & \texttt{url} (notebook session URL) \\ \hline
\texttt{notebook\_get\_cell} & \texttt{url}, \texttt{cell\_index} (1-based) \\ \hline
\texttt{insert\_cell} & \seqsplit{url, tab\_id, position, cell\_type, source} \\ \hline
\texttt{edit\_cell\_by\_index} & \seqsplit{url, tab\_id, cell\_index, source} \\ \hline
\texttt{delete\_by\_index} & \seqsplit{url, tab\_id, cell\_index} \\ \hline
\texttt{run\_cell} & \seqsplit{url, tab\_id, cell\_index} \\ \hline
\end{tabular}
\end{table}

The host coerces missing \texttt{url} and \texttt{tab\_id} fields from the active chat turn before dispatch, satisfying Objective~3 in Section~\ref{sec:objectives}.

\section{Host Configuration Flags}
\label{app:config}

Table~\ref{tab:app-config} lists environment variables governing Agentic behaviour during benchmarks.

\begin{table}[!ht]
\centering
\caption{Agentic configuration flags (\texttt{config.py}).}
\label{tab:app-config}
\footnotesize
\setlength{\tabcolsep}{3pt}
\renewcommand{\arraystretch}{1.02}
\begin{tabular}[t]{|>{\raggedright\arraybackslash\ttfamily\footnotesize}p{3.85cm}|>{\raggedright\arraybackslash}p{0.65cm}|>{\raggedright\arraybackslash}p{6.3cm}|}
\hline
\textbf{Variable} & \textbf{Default} & \textbf{Effect} \\ \hline
\texttt{LLM\_AGENTIC\_ENABLED} & on & Enables Agentic tool schemas \\ \hline
\texttt{AGENTIC\_MANDATORY\_}\newline\texttt{TWO\_PHASE} & on & Round~0 read-only; round~1 implement \\ \hline
\texttt{AGENTIC\_MAX\_TOOL\_ROUNDS} & 2 & Hard cap on LLM API calls per message \\ \hline
\texttt{AGENTIC\_FIRE\_AND\_FORGET} & on & Dispatch batch without per-cell replanning \\ \hline
\texttt{CEREBRAS\_MODEL} & \texttt{zai-glm-4.7} & GLM~4.7 model identifier \\ \hline
\end{tabular}
\end{table}

\section{Reproduction Commands}
\label{app:reproduce}

From the repository root, with \texttt{CEREBRAS\_API\_KEY} set in \texttt{testing/host/.env}:

\begin{verbatim}
python testing/host/scripts/run_agent_tests.py --live-llm
python testing/host/scripts/run_ask_tests.py --live-llm
python testing/host/scripts/run_code_tests.py --live-llm
python testing/host/scripts/generate_fyp_dissertation_report.py
\end{verbatim}

Log outputs appear under \texttt{testing/host/data/logs/}, including per-test JSON results and trace files ({\ttfamily\seqsplit{agentic\_tool\_trace.jsonl}}). The consolidated report is written to {\ttfamily\seqsplit{FYP\_DISSERTATION\_BENCHMARK\_REPORT.md}}.

\section{Primary Agentic Trace Summary}
\label{app:trace}

For Test~1 (ML Pipeline), the harness recorded the following dispatch profile in the second \ac{API} call:

\begin{itemize}
  \item Round~0: \texttt{notebook\_list\_cells}, \texttt{notebook\_get\_cell} (read phase);
  \item Round~1: 24$\times$ \texttt{insert\_cell}, 24$\times$ \texttt{edit\_cell\_by\_index}, 24$\times$ \texttt{run\_cell};
  \item Total tools in trace: 37 (including reads);
  \item Wall time: approximately 9.8\,s;
  \item Repair rounds: 0.
\end{itemize}

Host-side batch reordering applied descending index order for bulk inserts before run operations, consistent with Section~\ref{sec:tool-registry}.

\section{Extension Source Modules}
\label{app:extension}

Table~\ref{tab:app-ext} maps primary extension JavaScript modules to responsibilities.

\begin{table}[!ht]
\centering
\caption{Chrome extension modules.}
\label{tab:app-ext}
\small
\begin{tabular}{|p{4.2cm}|p{8.4cm}|}
\hline
\textbf{Module} & \textbf{Responsibility} \\ \hline
\texttt{background.js} & Native port, tab routing, scrape scheduling \\ \hline
\texttt{ui\_injector.js} & Chat panel injection on Kaggle pages \\ \hline
\texttt{cell\_index\_utils.js} & Positional index normalisation \\ \hline
\texttt{kernel\_state\_listener.js} & Execution state observation \\ \hline
\texttt{prompt\_observer.js} & User prompt capture for debugging \\ \hline
\end{tabular}
\end{table}

\section{Host Python Modules}
\label{app:host-modules}

Table~\ref{tab:app-host} lists core host modules referenced in Chapter~\ref{c-methods}.

\begin{table}[!ht]
\centering
\caption{Python native host modules.}
\label{tab:app-host}
\small
\begin{tabular}{|p{4.4cm}|p{8.2cm}|}
\hline
\textbf{Module} & \textbf{Responsibility} \\ \hline
\texttt{host.py} & Entry point; stdin/stdout native messaging loop \\ \hline
\texttt{dispatcher.py} & Message type routing \\ \hline
\texttt{streaming.py} & LLM streaming; mode resolution; tool rounds \\ \hline
\texttt{tool\_registry.py} & Tool registration, validation, dispatch \\ \hline
\texttt{agentic\_batch\_executor.py} & Batch queue build and reorder \\ \hline
\texttt{notebook\_data\_handler.py} & Ingest extension scrape payloads \\ \hline
\texttt{prompt\_engineering.py} & Message list assembly and context pack \\ \hline
\texttt{cerebras\_client.py} & HTTPS completions with tool schemas \\ \hline
\end{tabular}
\end{table}

\section{Glossary of Benchmark Identifiers}
\label{app:ids}

\begin{table}[!ht]
\centering
\caption{Benchmark test identifiers.}
\label{tab:app-ids}
\small
\begin{tabular}{|l|l|p{6.8cm}|}
\hline
\textbf{ID} & \textbf{Mode} & \textbf{Description} \\ \hline
ASK\_TEST\_1 & Ask & Notebook workflow explanation \\ \hline
ASK\_TEST\_2 & Ask & Cell~17 explanation \\ \hline
ASK\_TEST\_3 & Ask & Cell~31 error diagnosis \\ \hline
CODE\_TEST\_1 & Code & XGBoost pipeline \\ \hline
CODE\_TEST\_2 & Code & Cell~21 replacement \\ \hline
CODE\_TEST\_3 & Code & Feature-engineering module \\ \hline
CODE\_TEST\_4 & Code & Empty cell~50 deferral \\ \hline
AGENT\_TEST\_1 & Agentic & Complete ML pipeline (primary) \\ \hline
\end{tabular}
\end{table}

\section{Prompt Engineering Files}
\label{app:prompts}

Mode behaviour is controlled by text files under \texttt{testing/host/prompts/}. Table~\ref{tab:app-prompt-files} lists the principal templates.

\begin{table}[!ht]
\centering
\caption{Host prompt template files.}
\label{tab:app-prompt-files}
\footnotesize
\setlength{\tabcolsep}{3pt}
\renewcommand{\arraystretch}{1.02}
\begin{tabular}{|>{\raggedright\arraybackslash}p{4.45cm}|>{\raggedright\arraybackslash}p{7.25cm}|}
\hline
\textbf{File} & \textbf{Purpose} \\ \hline
\texttt{base\_notebook\_assistant.txt} & Core system persona and safety rules \\ \hline
\texttt{jupyter\_structure.txt} & Cell index and markdown/code conventions \\ \hline
\texttt{ask.txt} & Ask mode: explanation only, no tools \\ \hline
\texttt{code.txt} & Code mode: Placement + python block format \\ \hline
\texttt{agentic.txt} & Agentic mode: two-phase query--implement \\ \hline
\texttt{tool\_calling/react\_agent.txt} & Tool-calling behaviour for Agentic rounds \\ \hline
{\ttfamily\seqsplit{tool\_calling/tool\_descriptions\_autogen.txt}} & Auto-generated JSON schemas for tools \\ \hline
\end{tabular}
\end{table}

The Agentic template instructs the model to complete inspection in round~0 before emitting write tools in round~1. Code mode requires a \texttt{Placement} section before any fenced code block. Ask mode explicitly forbids code generation unless quoting existing notebook source for explanation.
\vspace{-6pt}

\section{Sample Ask Response Characteristics}
\label{app:sample-ask}

Although full model transcripts are omitted for brevity, reported Ask tests exhibited the following response structures:

\begin{itemize}
  \item \textbf{Test~1.} Sectioned prose (loading $\rightarrow$ cleaning $\rightarrow$ features $\rightarrow$ modelling $\rightarrow$ evaluation) with variable names matching the Pakistan housing snapshot.
  \item \textbf{Test~2.} Numbered list addressing inputs, outputs, dependencies, purpose, and risks for cell~17.
  \item \textbf{Test~3.} Root-cause paragraph naming \texttt{NameError}, minimal fix (restore \texttt{df} from upstream), and impact statement---without a full replacement cell.
\end{itemize}

These structures align with the rubric keywords defined in {\ttfamily\seqsplit{fyp\_experiment\_benchmarks\_ask.json}}.
\vspace{-4pt}

\section{Environment Setup}
\label{app:env}

Minimum environment requirements for reproduction:

\begin{itemize}
\setlength{\itemsep}{2pt}
  \item Google Chrome with the project extension loaded (developer mode).
  \item Python~3.10+ with dependencies from the host \texttt{requirements} specification.
  \item Registered native messaging host (\texttt{register.ps1} on Windows).
  \item Valid \texttt{CEREBRAS\_API\_KEY} in \texttt{testing/host/.env}.
\pagebreak[4]
  \item Frozen snapshots present under \texttt{data/notebooks/persistent/}.
\end{itemize}

Benchmark runners accept \texttt{--live-llm} to call the remote API; omitting this flag exercises mock paths suitable for continuous integration without API cost.

\section{Document Revision History}
\label{app:revision}

\begin{table}[!ht]
\centering
\caption{Thesis document revision summary.}
\label{tab:app-revision}
\small
\begin{tabular}{|l|l|p{7.8cm}|}
\hline
\textbf{Date} & \textbf{Version} & \textbf{Changes} \\ \hline
June 2026 & 1.0 & Initial submission: Chapters~1--7, benchmark appendix, APA references \\ \hline
June 2026 & 1.1 & Expanded evaluation chapters; dispatch-only Agentic framing; 85-page target \\ \hline
\end{tabular}
\end{table}

This revision aligns reported Agentic evidence with Test~1 (ML pipeline) only, documents Ask Tests~1--3 and Code Tests~1--4 per the approved evaluation protocol, and provides full reproducibility material for the FYP defence.
